\subsection{可交换元素。交换律}

\begin{enumerate}
\def\labelenumi{(\arabic{enumi})}
\setcounter{enumi}{3}
\tightlist
\item
  设 \((x, y) \mapsto x \top y\) 是 E 上的一个交换律；该律
\end{enumerate}

\[
(\mathbf {X}, \mathbf {Y}) \mapsto \mathbf {X} \top \mathbf {Y}
\]

在 E 的子集之间是交换的。

定义 9. 设 E 是一个广群，X 是 E 的一个子集。E 中与 X 的每个元素都可交换的元素所组成的集合称为 X 的中心化子。

设 X 和 Y 是 E 的两个子集，X' 和 Y' 分别是它们的中心化子。若 X ⊂ Y，则 Y' ⊂ X'。

设 \((\mathbf{X}_i)_{i\in I}\) 是 \(E\) 的一个子集族，且对所有 \(i\in I\) ，设 \(\mathbf{X}_i'\) 为 \(\mathbf{X}_i\) 的中心化子。 \(\bigcup_{i\in I}\mathbf{X}_i\) 的中心化子是 \(\bigcap_{i\in I}\mathbf{X}_i'\) .

设X是E的子集， \(X'\) 为X的中心化子。 \(X''\) （ \(X'\) 的中心化子）称为X的双中心化子。则 \(X \subset X''\) 。 \(X'''\) （ \(X''\) 的中心化子）等于 \(X'\) 。因为 \(X'\) 包含于它的双中心化子 \(X'''\) 中，且关系 \(X \subset X''\) 蕴含 \(X''' \subset X'\) .

命题3. 设E是一个结合广群，其律记为 \(\top\) 。如果一个元素 \(x\) （属于 E ）与元素 \(y\) 和 \(z\) （都属于 E ）可交换，则它与 \(y \top z\) 可交换.

对于

\[
\begin{array}{r l} x \top (y \top z) = (x \top y) \top z & = (y \top x) \top z \\ & = y \top (x \top z) = y \top (z \top x) = (y \top z) \top x. \end{array}
\]

推论。设E是一个结合广群。E的任意子集的中心化子是E的一个稳定子集。

定义10. 一个广群E的中心化子称为E的中心。E的中心中的元素称为E的中心元素。

如果E是一个结合广群，则根据命题3的推论，其中心是一个稳定子集，且在其中心上诱导的律是交换的。

命题4。设E为一个结合广群，X和Y为E的两个子集。如果X中的每个元素都与Y中的每个元素交换，那么由X生成的稳定子集中的每个元素都与由Y生成的稳定子集中的每个元素交换。

设 \(X'\) 和 \(X''\) 分别为X的中心化子和双中心化子。它们是E的稳定子集。现在 \(X \subset X''\) 且 \(Y \subset X'\) ，因此 \(X''\) （相应地， \(X'\) ）包含由 X（相应地，Y）生成的 E 的稳定子集。由于 \(X''\) 的每个元素与 \(X'\) 的每个元素可交换，命题得证。

推论1. 若 \(x\) 和 \(y\) 在结合律 \(\top\) 下可交换，则 \(\top^m x\) 和 \(\top^n y\) 对所有整数 \(m > 0\) 和 \(n > 0\) 亦然；特别地， \(\top^m x\) 和 \(\top^n x\) 对所有 \(x\) 和所有整数 \(m > 0\) 和 \(n > 0\) 都可交换.

推论2. 若子集X中任意两个元素在结合律T下可交换，则T在由X生成的稳定子集上诱导的律满足结合律与交换律。

定理2（交换性定理）。设 \(\top\) 是 E 上的一个结合合成律；设 \((x_{\alpha})_{\alpha \in A}\) 是 E 的元素的一个非空有限族，其元素两两可交换；设B和C是两个以A为底集的全序集。则 \(\bigcap_{\alpha \in B} x_{\alpha} = \bigcap_{\alpha \in C} x_{\alpha}\) .

由于当A只有一个元素时定理成立，我们对A中元素个数p进行归纳论证。设p为整数 \textgreater1，并假设当Card A \textless{} p.～我们证明当Card A = p时成立。可以假设 A 是 N 中的区间 \([0, p - 1]\) ；有序序列 \((x_{\alpha})_{\alpha \in \mathbf{A}}\) （由 \(\mathbf{A}\) 上的自然序关系定义）的合成是 \(\prod_{i = 0}^{p - 1}x_i\) .

设 A 被赋予另一个全序，并设 h 是该序下 A 的最小元素，且 \(A'\) 为 A 的其余元素所成的集合（按诱导序全序）。首先假设 0 \textless{} h \textless{} p - 1 ，令 \(P = \{0, 1, \ldots, h - 1\}\) 和 \(Q = \{h + 1, \ldots, p - 1\}\) ；定理对 \(A'\) 已假设成立，应用结合性定理，我们得到（因为 \(A' = P \cup Q\) )

\[
\bigcap_ {\alpha \in \mathbf {A} ^ {\prime}} x _ {\alpha} = \left(\prod_ {i = 0} ^ {h - 1} x _ {i}\right) \top \left(\prod_ {i = h + 1} ^ {p - 1} x _ {i}\right)
\]

由此，将 \(x_{h}\) 与两边复合，并反复应用 T 的交换律和结合律：

\[
\begin{array}{r l} \underset {\alpha \in \mathbf {A}} {\top} x _ {\alpha} & = x _ {h} \top \left(\underset {\alpha \in \mathbf {A} ^ {\prime}} {\top} x _ {\alpha}\right) = x _ {h} \top \left(\underset {i = 0} {\overset {h - 1} {\top}} x _ {i}\right) \top \left(\underset {i = h + 1} {\overset {p - 1} {\top}} x _ {i}\right) \\ & = \left(\underset {i = 0} {\overset {h - 1} {\top}} x _ {i}\right) \top x _ {h} \top \left(\underset {i = h + 1} {\overset {p - 1} {\top}} x _ {i}\right) = \underset {i = 0} {\overset {p - 1} {\top}} x _ {i}, \end{array}
\]

这便证明了该情形下的定理。若 h = 0 或 h = p - 1，同样的结论成立，但更为简单，因为由 P 产生的项或由 Q 产生的项不会出现在公式中。

在集合 E 上的一个交换结合律下，有限族 \((x_{\alpha})_{\alpha \in A}\) （由 E 的元素组成）的合成按定义是：将 A 以一切可能的方式全序化后所得的所有有序序列的合成的公共值。这个合成在由 T 表示的律下仍记作 \(\bigcap_{\alpha \in A} x_{\alpha}\) ；其他记号也类似。

定理 3. 设 \(\top\) 是 E 上的一个结合律，且 \((x_{\alpha})_{\alpha \in \mathbf{A}}\) 为 E 的元素的一个非空有限族，其元素两两可交换。若 A 是非空子集 \((\mathbf{B}_i)_{i\in I}\) （两两不相交）的并集，则

\[
\bigcap_ {\alpha \in A} x _ {\alpha} = \bigcap_ {i \in I} \left(\bigcap_ {\alpha \in B _ {i}} x _ {\alpha}\right).\tag{6}
\]

这由定理2推出，若 A 和 I 已全序化，使得诸 \(B_{i}\) 满足定理1的条件。

我们特别指出该定理的两个重要特例：

\begin{enumerate}
\def\labelenumi{\arabic{enumi}.}
\tightlist
\item
  若 \((x_{\alpha\beta})_{(\alpha,\beta)\in\mathbf{A}\times\mathbf{B}}\) 是结合广群中可交换元素的一个有限族，其指标集是两个非空有限集 A、B 的乘积（一个“双重族”），则
\end{enumerate}

\[
\underset {(\alpha , \beta) \in A \times B} {\top} x _ {\alpha \beta} = \underset {\alpha \in A} {\top} \left(\underset {\beta \in B} {\top} x _ {\alpha \beta}\right) = \underset {\beta \in B} {\top} \left(\underset {\alpha \in A} {\top} x _ {\alpha \beta}\right)\tag{7}
\]

这由定理3推出：一方面把 \(\mathbf{A} \times \mathbf{B}\) 看作诸集合 \(\{\alpha\} \times \mathbf{B}\) 的并集，另一方面看作诸集合 \(\mathbf{A} \times \{\beta\}\) 的并集。

例如，若B有n个元素，且对每个 \(\alpha \in A\) 所有 \(x_{\alpha\beta}\) 具有相同的值 \(x_{\alpha}\) ，则

\[
\bigcap_ {\alpha \in \mathbf {A}} \left(\frac {n}{T} x _ {\alpha}\right) = \frac {n}{T} \left(\bigcap_ {\alpha \in \mathbf {A}} x _ {\alpha}\right).\tag{8}
\]

如果B有两个元素，我们得到以下结果：设 \((x_{\alpha})_{\alpha \in A}, (y_{\alpha})_{\alpha \in A}\) 为E的两个非空元素族。如果这些 \(x_{\alpha}\) 和 \(y_{\beta}\) 两两可交换，则

\[
\underset {\alpha \in \mathbf {A}} {\top} x _ {\alpha} \top y _ {\alpha} = \left(\underset {\alpha \in \mathbf {A}} {\top} x _ {\alpha}\right) \top \left(\underset {\alpha \in \mathbf {A}} {\top} y _ {\alpha}\right).\tag{9}
\]

由于公式(7)，双序列 \((x_{ij})\) （其指标集为两个区间 \(\{p, q\}\) 与 \(\{r, s\}\) （在 \(\mathbf{N}\) 中）的乘积）的合成，对于以加法记法书写的交换结合律，通常记为

\[
\sum_ {i = p} ^ {q} \sum_ {j = r} ^ {s} x _ {i j} \quad \text { or } \quad \sum_ {j = r} ^ {s} \sum_ {i = p} ^ {q} x _ {i j}
\]

对于用其他符号表示的律亦然。

\begin{enumerate}
\def\labelenumi{\arabic{enumi}.}
\setcounter{enumi}{1}
\tightlist
\item
  设 \(n\) 为整数 \(>0\) ，并设 \(A\) 为整数有序对 \((i,j)\) （满足 \(0 \leqslant i \leqslant n, 0 \leqslant j \leqslant n\) 和 \(i < j\) ）所成的集合；族 \((x_{ij})_{(i,j) \in A}\) 的合成（在交换结合律下）也记为 \(\bigcap_{0 \leqslant i < j \leqslant n} x_{ij}\) （或简记为 \(\bigcap_{i < j} x_{ij}\) （若无混淆之虞）；此处定理3给出公式
\end{enumerate}

\[
\underset {0 \leqslant i <   j \leqslant n} {\mathsf {T}} x _ {i j} = \underset {i = 0} {\overset {n - 1} {\mathsf {T}}} \left(\underset {j = i + 1} {\overset {n} {\mathsf {T}}} x _ {i j}\right) = \underset {j = 1} {\overset {n} {\mathsf {T}}} \left(\underset {i = 0} {\overset {j - 1} {\mathsf {T}}} x _ {i j}\right).\tag{10}
\]

对于指标集为多于两个集合之积的族，存在与(7)类似的公式；对于指标集为集合 \(S_{p}\) （由严格递增序列 \((i_{k})_{1\leqslant k\leqslant p}\) （p 个整数，满足 \(0\leqslant i_{k}\leqslant n\) ( \(p\leqslant n+1\) )）所组成）的族，存在与(10)类似的公式：在后一种情形下，族 \((x_{i_{1}i_{2}\ldots i_{p}})(i_{1},\ldots,i_{p})\in\mathbb{S}_{p}\) 的合成记为

\[
\underset {0 \leqslant i _ {1} <   i _ {2} <   \dots <   i _ {p} <   n} {\mathsf {T}} x _ {i _ {1} i _ {2} \ldots i _ {p}}, \quad \text {or simply} \quad \underset {i _ {1} <   i _ {2} <   \dots <   i _ {p}} {\mathsf {T}} x _ {i _ {1} i _ {2} \ldots i _ {p}}.
\]

命题5. 设 E 和 F 为两个广群，其律都记为 \(\top\) ，并设 \(f\) 和 \(g\) 为从 E 到 F 的同态。设 \(f \top g\) 为映射 \(x \mapsto f(x) \top g(x)\) （从 E 到 F ）。若 F 是结合的且交换的，则 \(f \top g\) 是一个同态。

对所有元素 \(x\) 和 \(y\) （属于 \(\mathbf{E}\) ）：

\[
\begin{array}{r l} (f \top g) (x \top y) & = f (x \top y) \top g (x \top y) = f (x) \top f (y) \top g (x) \top g (y) \\ & = f (x) \top g (x) \top f (y) \top g (y) = ((f \top g) (x)) \top ((f \top g) (y)). \end{array}
\]

